% !TEX root = ../../main.tex

\chapter{神经网络训练与泛化}

我们已经学习了反向传播算法，并进一步掌握了学习率调度、动量法、自适应优化器等一系列参数更新策略。现在，我们构建一个神经网络并尝试训练它，看看训练过程中会发生什么。

\begin{figure}[H]
    \boxedimage[width=0.95\textwidth, trim=0 0 0 0, clip]{figures/part_01/p1_ch04/network.pdf}
    \caption{10层神经网络的结构示意图。输入$x$依次经过9个隐藏层$h^{(1)}, h^{(2)}, \dots, h^{(9)}$的逐层变换，最终得到输出$y$。}
    \label{fig:network}
\end{figure}

假设我们构建一个如\cref{fig:network}所示的10层神经网络，其中包含9层隐藏层 $h^{(1)} \sim h^{(9)}$和1个输出层$y$（输入不计入层数），并且每一层都使用Sigmoid作为激活函数。在前向传播计算得到预测值与损失函数$\mathcal{L}$后，我们通过反向传播计算各层参数的梯度，并据此更新参数。我们以第一层参数$w^{(1)}$的梯度为例。根据链式法则，损失函数对$w^{(1)}$的梯度需要从输出层一路向前传递：
$$
\frac{\partial \mathcal{L}}{\partial w^{(1)}} = \frac{\partial \mathcal{L}}{\partial y} \cdot \frac{\partial y}{\partial h^{(9)}} \cdot \frac{\partial h^{(9)}}{\partial h^{(8)}} \cdots \frac{\partial h^{(2)}}{\partial h^{(1)}} \cdot \frac{\partial h^{(1)}}{\partial w^{(1)}}
$$
可以看到，第一层的梯度并不是由某一个单独的导数决定的，而是一长串局部导数的连乘。对于网络中的每一层，反向传播时的局部梯度同时受到权重和激活函数导数的影响。为了直观理解，我们暂时把每层看作一个简单的标量变换（记$h^{(0)} = x$，即输入）：
$$
h^{(l)}=\sigma\left(w^{(l)}h^{(l-1)}\right)
$$
那么这一层在反向传播过程中产生的梯度因子就是：
$$
\frac{\partial h^{(l)}}{\partial h^{(l-1)}}=w^{(l)}\sigma'\left(w^{(l)}h^{(l-1)}\right)
$$
因此，随着网络层数增加，权重$w^{(l)}$和激活函数导数$\sigma'$都会参与连续相乘，它们共同决定梯度究竟是逐渐变小，还是不断变大。

\begin{figure}[H]
    \boxedimage[width=0.8\textwidth]{figures/part_01/p1_ch02/sigmoid.pdf}
    \caption{Sigmoid函数与导数。左图$f(x)$在$x=0$处取$0.5$、两端饱和；右图$f'(x)$在$x=0$处峰值$0.25$。}
    \label{fig:sigmoid_function}
\end{figure}

那么，Sigmoid的导数表现如何呢？如\cref{fig:sigmoid_function}所示，Sigmoid导数的最大值仅为$0.25$；而当输入逐渐远离$0$时，Sigmoid会进入饱和区，导数也会迅速接近$0$。

现在，我们先考虑一个最简单的情况：假设每层的权重都初始化为$1$，并且每层的Sigmoid导数都恰好取到最大值$0.25$。此时，每层的梯度因子就是
$$
w^{(l)}\sigma'\left(w^{(l)}h^{(l-1)}\right) = 1\times0.25 = 0.25
$$
经过10层连续相乘后，梯度因子变为
$$
0.25^{10} \approx 9.5\times10^{-7}
$$
也就是说，即使每一层都处于最有利的状态，梯度经过10层反向传播后，也已经衰减到了不到百万分之一。而在实际网络中，Sigmoid并不总是处于$x=0$附近，许多神经元可能进入饱和区，此时导数会远小于$0.25$，梯度衰减自然会更加严重。随着网络不断加深，传递到前面层的梯度可能越来越小，小到几乎无法有效更新参数，这就是\textbf{梯度消失（Gradient Vanishing）}。

而如果权重的尺度过大，情况就可能完全相反。假设某一层的权重为$20$，同时该层的激活函数仍处于$x=0$附近，使得$\sigma'(x)\approx0.25$，那么这一层的梯度因子约为
$$
w^{(l)}\sigma'\left(w^{(l)}h^{(l-1)}\right) \approx 20\times0.25 = 5
$$
如果连续10层都具有类似的梯度因子，那么经过10层连乘后，梯度将被放大为：
$$
5^{10}=9,765,625
$$
也就是说，梯度在反向传播过程中被放大了近千万倍。随着网络进一步加深，这种放大效应会更加严重，使参数更新变得剧烈甚至失控，严重时还可能出现数值溢出，最终导致Loss变成\texttt{NaN}。这就是\textbf{梯度爆炸（Gradient Exploding）}。

可以看到，梯度消失和梯度爆炸本质上都来自反向传播中的连乘。每一层的梯度因子都由权重和激活函数导数共同决定：当它们的整体尺度持续小于$1$时，梯度会随着网络深度增加而不断衰减；当整体尺度持续大于$1$时，梯度则会迅速膨胀。这也说明，深度网络的训练不仅依赖反向传播和优化算法，还面临着梯度传播、训练稳定性以及模型泛化等一系列问题。接下来，我们将从梯度消失和梯度爆炸开始，进一步讨论这些问题及其相应的解决方法。

\section{梯度消失}

既然Sigmoid的导数最大只有$0.25$，梯度会随着网络加深而不断衰减，那我们是不是可以寻找一个导数始终保持在$1$附近的激活函数，让梯度在反向传播过程中尽可能少地衰减？如\cref{fig:x}所示，最简单的选择就是$f(x)=x$，它在每一点的导数都恒为$1$。这样一来，在不考虑权重影响的情况下，梯度经过这一层时不会因为激活函数本身而产生额外的衰减，可以直接传递到前一层。

\begin{figure}[H]
    \boxedimage[width=0.8\textwidth]{figures/part_01/p1_ch04/x.pdf}
    \caption{线性函数$f(x)=x$及其导数。左图为过原点、斜率为1的直线；右图其导数$f'(x)$恒为1。}
    \label{fig:x}
\end{figure}

但是，我们之前已经学过，激活函数最核心的作用是引入非线性。而$f(x)=x$本身就是一个线性函数。如果每一层都使用它作为激活函数，那么无论网络堆叠多少层，多个线性变换最终都可以合并为一个线性变换。这样一来，深层网络就失去了建模复杂非线性关系的能力，甚至连最简单的异或问题都无法解决。

因此，我们需要找到一种激活函数，在保留一定非线性的同时，又能够让梯度更加稳定地传播。一个直接的想法是，我们可以把$f(x)=x$“折断”：保留正半轴的部分，而将负半轴直接折断为$0$。于是，我们可以得到：
$$
f(x) = \max(0, x)
$$

\begin{figure}[H]
    \boxedimage[width=0.8\textwidth]{figures/part_01/p1_ch04/relu.pdf}
    \caption{ReLU函数及其导数。左图$f(x)$在$x<0$时恒为0、$x\geq 0$时为斜率1的直线；右图其导数$f'(x)$在$x<0$时为0、$x>0$时恒为1，在$x=0$处不可导。}
    \label{fig:relu}
\end{figure}

如\cref{fig:relu}所示，这就是\textbf{ReLU（Rectified Linear Unit）函数}\footnote{ReLU最早由Vinod Nair和Geoffrey Hinton于2010年提出，论文为\href{https://www.cs.toronto.edu/~fritz/absps/reluICML.pdf}{\textit{Rectified Linear Units Improve Restricted Boltzmann Machines}}。}。它的形式非常简单，却很好地兼顾了非线性表达能力和梯度传播效果。首先，当$x>0$时，ReLU的导数恒为$1$。与Sigmoid相比，梯度在正区间不会因为激活函数本身而逐层缩小，因此能够显著缓解深层网络中的梯度消失问题。其次，ReLU只需要进行简单的比较和取值，不需要计算Sigmoid中的指数函数，因此计算更加高效。最后，当$x<0$时，ReLU的输出恒为$0$，因此在一次前向传播中只有部分神经元处于激活状态，从而形成稀疏激活。这种选择性的响应方式能够减少冗余表示，使不同神经元更容易学习不同的特征。

当然，ReLU也并非没有问题。最典型的问题就是\textbf{Dead ReLU}。当$x<0$时，ReLU的导数恒为$0$。如果某个神经元在训练过程中进入负半轴，并且之后对输入产生的结果始终小于$0$，那么该神经元的梯度就会持续为$0$，参数也无法通过梯度下降得到有效更新，最终可能使这个神经元长期处于“死亡”状态。尽管如此，ReLU凭借简单的计算形式、良好的非线性表达能力以及在正区间稳定的梯度传播特性，成为深度神经网络中最经典、应用最广泛的激活函数之一。

但需要注意的是，ReLU主要解决的是连乘因子中激活函数导数这一部分的问题，但别忘了，链式法则中的每一项$\frac{\partial h^{(l)}}{\partial h^{(l-1)}}$实际上是激活函数导数$\sigma'(z)$与权重$w$的乘积。激活函数的问题解决了，但权重仍然可能成为梯度传播的瓶颈：如果权重初始化得太小，连乘因子的整体尺度仍然可能小于$1$，梯度依然会逐层衰减；如果权重初始化得太大，连乘因子的整体尺度则可能大于$1$，进而导致梯度爆炸。因此，在选择合适的激活函数之后，我们还需要进一步考虑权重应当从什么样的状态开始训练，使网络在训练初期就处于合适的尺度。这就是\textbf{权重初始化（Weight Initialization）}。

那我们能不能把权重全部初始化为$0$，或者全部初始化为同一个数值呢？我们假设同一层中的所有神经元都使用完全相同的权重，那么它们接收到相同的输入后，就会计算出完全相同的输出；在反向传播时，它们又会得到完全相同的梯度。于是，每次参数更新之后，它们仍然保持完全一致。无论训练多少轮，这些神经元学到的内容都相同，最终相当于只有一个神经元在工作，网络也就失去了利用多个神经元学习不同特征的意义。

因此，我们需要在初始化时引入随机性，让不同神经元从不同的参数起点出发，从而打破这种对称性。但打破对称性只是第一步，权重也不能“随便随机”。如果权重随机得太小，那么前向传播时信号会逐层变弱，反向传播时梯度也可能不断衰减；反过来，如果权重随机得太大，信号和梯度又可能在层层传递中不断放大，最终导致训练不稳定。也就是说，我们不仅需要让不同神经元从不同的地方出发，还需要让权重的初始尺度保持在一个合适的范围内。

那么，什么样的随机初始化才是合适的？我们先从前向传播的角度进行简单推导。考虑某一层的加权求和：
$$
z=\sum_{i=1}^{n_{in}}w_ix_i
$$
假设权重与输入相互独立，并且二者的均值都为$0$，根据独立随机变量和的方差性质，则有：
$$
\operatorname{Var}(z) = n_{in}\operatorname{Var}(w)\operatorname{Var}(x)
$$
如果希望信号经过这一层后，其方差保持在大致相同的尺度，即
$$
\operatorname{Var}(z)\approx\operatorname{Var}(x)
$$
那么就需要满足
$$
\operatorname{Var}(w)\approx\frac{1}{n_{in}}
$$
这是从前向传播得到的要求。同样地，如果从反向传播的角度出发，希望梯度在逐层传递时也保持相近的尺度，则会得到与$n_{out}$相关的另一个要求：
$$
\operatorname{Var}(w)\approx\frac{1}{n_{out}}.
$$
因此，为了同时兼顾前向传播和反向传播，可以在二者之间取一个平衡，将权重的方差设置为：
$$
\operatorname{Var}(w)=\frac{2}{n_{in}+n_{out}}
$$
这种初始化方法就是\textbf{Xavier初始化}\footnote{Xavier初始化由Xavier Glorot和Yoshua Bengio于2010年提出，原始论文为\href{https://proceedings.mlr.press/v9/glorot10a/glorot10a.pdf?ref=blog.kusho.ai}{\textit{Understanding the Difficulty of Training Deep Feedforward Neural Networks}}。}。它主要针对Sigmoid、Tanh这类在$0$附近近似线性、输出以$0$为中心的激活函数设计，通过控制权重的方差，使信号和梯度在网络中传播时尽可能保持稳定的尺度。

\begin{orangeBox}{Tanh激活函数}
Tanh可以看作Sigmoid的一个中心化版本。Sigmoid的输出范围是$(0,1)$，而Tanh则将其重新映射到了$(-1,1)$，使其以$0$为中心。

\begin{figure}[H]
    \boxedimage[width=0.8\textwidth]{figures/part_01/p1_ch04/tanh.pdf}
    \caption{Tanh函数及其导数。左图$f(x)$在$x=0$处取0、两端趋于$\pm 1$并饱和；右图其导数$f'(x)$在$x=0$处峰值1、随$|x|$增大衰减至零。}
    \label{fig:tanh}
\end{figure}

Tanh的数学形式为：
$$
\tanh(x)=\frac{e^x-e^{-x}}{e^x+e^{-x}}.
$$
其输出范围为$(-1,1)$，并且满足$\tanh(0)=0$。其导数为
$$
\tanh'(x)=1-\tanh^2(x).
$$
在$x=0$处，Tanh的导数最大，为$1$；当输入逐渐远离$0$时，函数进入饱和区，导数也会逐渐接近$0$。因此，相比Sigmoid，Sigmoid具有更好的零中心特性，但仍然属于饱和型激活函数，在深层网络中依然可能出现梯度消失问题。
\end{orangeBox}

需要注意的是，上面的推导得到的是权重方差的目标值，而在实际初始化时，还需要进一步确定权重从什么分布中采样。常见的实现方式有均匀分布和正态分布两种。对于均匀分布，假设
$$
w\sim\mathcal{U}(-a,a)
$$
其方差为：
$$
\operatorname{Var}(w)=\frac{a^2}{3}
$$
我们令其等于Xavier初始化要求的目标方差，则有：
$$
\frac{a^2}{3}=\frac{2}{n_{in}+n_{out}}
$$
从而可以得到：
$$
a=\sqrt{\frac{6}{n_{in}+n_{out}}}
$$
因此，Xavier均匀初始化可以写成：
$$
w\sim\mathcal{U}\left(-\sqrt{\frac{6}{n_{in}+n_{out}}},\sqrt{\frac{6}{n_{in}+n_{out}}}\right)
$$

而对于正态分布，我们则令权重服从均值为$0$、方差为$\frac{2}{n_{in}+n_{out}}$的高斯分布：
$$
w\sim\mathcal{N}\left(0,\frac{2}{n_{in}+n_{out}}\right)
$$
此时对应的标准差为：
$$
\sigma=\sqrt{\frac{2}{n_{in}+n_{out}}}.
$$
因此，Xavier初始化并不是简单地随机生成一组权重，而是先确定合适的方差，再根据具体的概率分布进行采样，使网络在训练开始时就具有较为稳定的信号和梯度尺度。

然而，当激活函数换成ReLU时，情况又发生了变化。前面我们已经看到，ReLU会将所有负值直接截断为$0$。在输入近似关于$0$对称的情况下，大约有一半的值会被置为$0$，从而使信号的尺度受到压缩。如果继续使用Xavier初始化，随着网络不断加深，信号可能逐层减弱，重新出现梯度衰减的问题。针对ReLU的这一特点，我们可以进一步推导合适的权重尺度。设线性变换后的结果为：
$$
z=\sum_{i=1}^{n_{in}}w_ix_i
$$
在输入近似零均值且关于$0$对称的情况下，$z$也可以近似看作关于$0$对称。为了直观理解这一点，假设$z$取值为
$$
-2,-1,0,1,2
$$
经过ReLU处理后变为
$$
0,0,0,1,2
$$
可以看到，原来正负两侧都有取值，经过ReLU后，负半部分被全部截断为$0$。对于关于$0$对称的分布，正半部分和负半部分对$z^2$的贡献大致相同，而ReLU只保留正半部分，因此保留下来的平方值总量约为原来的一半。也就是：
$$
\mathbb{E}[\operatorname{ReLU}(z)^2]\approx\frac{1}{2}\mathbb{E}[z^2]
$$
如果希望经过这一层后信号的尺度保持稳定，就需要满足：
$$
\frac{1}{2}n_{in}\operatorname{Var}(w)\operatorname{Var}(x)\approx\operatorname{Var}(x)
$$
由此可以得到：
$$
\operatorname{Var}(w)=\frac{2}{n_{in}}.
$$
这里的系数$2$正是为了补偿ReLU截断负值带来的约一半的信号损失，这就是\textbf{He初始化（He Initialization）}的核心公式\footnote{He初始化由Kaiming He、Xiangyu Zhang、Shaoqing Ren和Jian Sun于2015年提出，原始论文为\href{https://www.cv-foundation.org/openaccess/content_iccv_2015/papers/He_Delving_Deep_into_ICCV_2015_paper.pdf}{\textit{Delving Deep into Rectifiers: Surpassing Human-Level Performance on ImageNet Classification}}。}。

在实际初始化时，同样可以选择不同的概率分布。对于均匀分布，有：
$$
w\sim\mathcal{U}\left(-\sqrt{\frac{6}{n_{in}}},\sqrt{\frac{6}{n_{in}}}\right)
$$
而对于正态分布，则有：
$$
w\sim\mathcal{N}\left(0,\frac{2}{n_{in}}\right).
$$
相比Xavier初始化，He初始化进一步增大了权重的方差，以补偿ReLU截断负值带来的尺度损失，因此更加适合ReLU及其变体，也成为现代深度神经网络中常用的参数初始化方法之一。

\section{梯度爆炸}

那么，梯度爆炸又该怎么解决呢？首先，前面介绍的参数初始化本身就可以在一定程度上缓解梯度爆炸。我们已经知道，梯度爆炸和梯度消失本质上都来自反向传播中的连乘，只是前者对应的连乘因子整体偏大，后者则整体偏小。因此，Xavier、He等初始化方法通过合理控制权重的尺度，使梯度因子在训练开始时尽可能保持在合适的范围内，从源头上降低梯度爆炸发生的概率。但初始化只能控制训练开始时的状态。在实际训练过程中，随着参数不断更新，梯度仍然可能因为随机采样、异常数据等因素突然变得很大。此时，我们还需要一种能够在训练过程中直接限制梯度大小的方法。

最直接的想法是调小学习率$\eta$，让参数更新$\eta\cdot g$不要过大。但这样会连正常的梯度一起缩小，导致整个训练过程变慢。而我们真正想限制的，是那些异常大的梯度，而不是所有梯度。因此，与其调整学习率，不如直接限制梯度本身。这就是\textbf{梯度裁剪（Gradient Clipping）}的核心思想\footnote{梯度裁剪由Razvan Pascanu、Tomas Mikolov和Yoshua Bengio于2013年提出，原始论文为\href{https://proceedings.mlr.press/v28/pascanu13.pdf}{\textit{On the Difficulty of Training Recurrent Neural Networks}}。}：在反向传播得到梯度之后、更新参数之前，先检查梯度的大小；如果梯度超过预设阈值，就将其按比例缩小到阈值以内。以按范数裁剪为例：

$$
g\leftarrow
\begin{cases}
    \dfrac{\text{threshold}}{\|g\|}g,&\|g\|>\text{threshold},\\
    g,&\text{otherwise}.
\end{cases}
$$
这种方法不会改变梯度的方向，只会限制梯度的幅度。例如，假设阈值设为$5$，而某次反向传播得到梯度$g=(30,40)$，其范数为：
$$
\|g\|=\sqrt{30^2+40^2}=50
$$
由于$50>5$，我们需要将梯度缩放为原来的$5/50=0.1$，可以得到：
$$
g=(3,4)
$$
裁剪前后两个分量的比例始终是$3:4$，方向完全没有改变，只是幅度被拉回了安全范围。因此，合理的参数初始化可以从训练起点缓解梯度爆炸，而梯度裁剪则可以在训练过程中限制异常大的梯度。

\section{批量归一化}

我们现在已经通过各种方法缓解了梯度消失和梯度爆炸，按理说，深层网络应该能够更加稳定地训练了。但事情往往并没有这么简单。

我们来看一个具体的例子。假设我们观察网络的第10层。在训练初期，这一层接收到的输入可能集中在均值$0.2$、标准差$0.4$附近；随着前面第9层参数不断更新，它的输出分布也会随之变化，第10层接收到的输入可能逐渐变成均值$0.8$、标准差$1.5$，甚至进一步变化到均值$6$、标准差$3$。这样一来，第10层就需要不断适应变化中的输入分布。尤其当使用Sigmoid或Tanh时，如果输入逐渐偏离$0$，就更容易进入饱和区，使梯度变小，进一步增加训练难度。换句话说，随着前面层的参数不断更新，后面层接收到的特征分布也会持续发生变化。网络越深，这种变化越容易逐层累积，使后面的层始终需要适应不断变化的输入。这种由网络内部参数更新导致的中间特征分布持续变化的现象，被称为\textbf{内部协变量偏移（Internal Covariate Shift）}。

既然网络内部的特征分布会不断变化，那么我们能不能在每一层内部对输入进行规范化，让它始终保持在一个较为稳定的尺度？这正是\textbf{批量归一化（Batch Normalization，BN）}的基本思想\footnote{Batch Normalization由Sergey Ioffe和Christian Szegedy于2015年提出，原始论文为\href{https://proceedings.mlr.press/v37/ioffe15.html}{\textit{Batch Normalization: Accelerating Deep Network Training by Reducing Internal Covariate Shift}}。}。对于一个Mini-batch，BN首先计算每个特征维度上的均值$\mu_B$和方差$\sigma_B^2$，然后对输入进行标准化：
$$
\hat{x}_i=\frac{x_i-\mu_B}{\sqrt{\sigma_B^2+\epsilon}}.
$$
其中，$\epsilon$是一个很小的常数，例如$10^{-5}$，用于避免分母为$0$。经过这一步之后，当前Mini-batch中的数据会被重新调整到较为稳定的中心和尺度。

但是，这里又会产生一个新的问题：如果我们把所有数据都强行标准化，那么网络是不是只能接收均值接近$0$、尺度接近$1$的输入？如果原来的分布恰好更有利于当前任务，这些信息岂不是被直接破坏了？

为了解决这个问题，BN在标准化之后又引入了两个可学习的参数——缩放参数$\gamma$和平移参数$\beta$：
$$
y_i=\gamma\hat{x}_i+\beta.
$$
这样，网络就可以根据任务需要重新调整标准化后的数据。如果标准化后的结果已经足够合适，模型可以学习到接近$\gamma=1,\beta=0$的参数；如果需要不同的尺度和偏移，也可以通过训练自动调整$\gamma$和$\beta$。

批量归一化的整个过程如\cref{fig:batch_normalization}所示。
\begin{figure}[H]
    \boxedimage[width=0.95\textwidth]{figures/part\_01/p1\_ch04/batch_normalization.pdf}
    \caption{批归一化示意图。输入$x_i$（左，均值$\mu_B$、标准差$\sigma_B$）先被标准化为零均值、单位方差的$\hat{x}_i$（中），再经可学习参数$\gamma$、$\beta$缩放平移为$y_i$（右）。}
    \label{fig:batch_normalization}
\end{figure}

因此，BN并不是简单地将数据固定在某一种分布，而是通过“先标准化、再进行可学习变换”的方式，在稳定特征尺度的同时保留网络自主调整特征分布的能力。这样既能减小层间输入尺度不断变化带来的训练困难，也不会限制模型的表达能力。此外，训练过程中基于当前Mini-batch计算均值$\mu_B$和方差$\sigma_B^2$会引入一定的随机波动，从而产生轻微的隐式正则化效果，在一定程度上帮助缓解过拟合。总体而言，BN通过归一化中间特征并引入可学习的缩放与平移，使深层网络中的特征尺度更加稳定，从而降低模型对权重初始化和超参数设置的敏感性，并改善训练过程的稳定性与效率。

\begin{purpleBox}{现代视角：BN为什么有助于优化？}
虽然BN在工程实践中取得了巨大成功，但学界对其作用机制的认识也在不断深入。最初，原作者将其效果归结为消除了“内部协变量偏移”；然而，Shibani Santurkar等人在论文\href{https://arxiv.org/abs/1805.11604}{\textit{How Does Batch Normalization Help Optimization?}}中对这一观点进行了重新审视。研究人员通过对照实验发现，加入BN的网络中间层分布依然存在显著漂移；更具说服力的是，当实验者刻意在BN之后人为注入剧烈变化的均值和方差噪声，即人为制造较强的内部协变量偏移时，网络的训练速度和最终性能仍然基本不受影响。这表明，消除内部协变量偏移既不是BN发挥作用的必要条件，也不是解释其效果的核心原因。

研究进一步发现，BN能够改善优化问题的几何性质，使损失函数及其梯度在训练过程中更加平滑、稳定。这样一来，梯度的变化更加可控，优化器也能够使用更大的学习率进行参数更新，从而更快、更稳定地找到较好的解。因此，从现代视角来看，BN的重要价值不仅在于调整中间特征的尺度，更在于改善深层网络的优化过程，让参数更新更加稳定、高效，这也是BN能够成为深度神经网络中经典组件的重要原因。
\end{purpleBox}

经过这些方法，深层网络的训练可以更加稳定，但训练过程中的问题并不只有梯度和优化，模型的泛化能力同样值得关注。

\section{过拟合与正则化}

在实际训练中，我们通常先利用\textbf{训练集（Training Set）}训练模型，并根据需要划分\textbf{验证集（Validation Set）}选择超参数、比较模型（验证集非必须，数据有限时也可以不单独划分）；当模型确定后，再使用完全没有参与训练的\textbf{测试集（Test Set）}进行最终评估，以判断模型在未见数据上的\textbf{泛化能力（Generalization）}。下面假设在同一个二分类任务上进行三组实验，训练流程基本相同，仅改变模型容量和训练时长，最终在训练集和测试集上得到的准确率如表~\ref{tab:fitting}所示。

\begin{table}[H]
    \centering
    \caption{三种拟合情况在训练集与测试集上的准确率}
    \label{tab:fitting}
    \begin{tabular}{ccc}
        \toprule
        \thead{} & \thead{训练集准确率} & \thead{测试集准确率} \\
        \midrule
        情况一 & 65.2\% & 60.1\% \\
        情况二 & 93.4\% & 91.7\% \\
        情况三 & 99.8\% & 61.3\% \\
        \bottomrule
    \end{tabular}
\end{table}

情况一中，模型在训练集上的准确率就不高，在测试集上同样偏低，这属于\textbf{欠拟合（Underfitting）}。模型容量不足或训练不充分，使其尚未学到数据中的基本规律，因此在训练集和测试集上的表现都较差。情况二中，训练集与测试集的准确率都较高，且二者差距较小，这属于\textbf{拟合良好（Good Fit）}。模型在测试集上也能保持较好的表现，说明模型学到了具有一定普适性的规律，具有较好的泛化能力。情况三则最具迷惑性：训练集准确率接近$100\%$，但测试集准确率却明显下降，这就是\textbf{过拟合（Overfitting）}。模型虽然能够很好地拟合训练数据，却把其中的样本细节、噪声甚至偶然规律也一并记忆下来，因此一旦面对新的数据，性能便会明显下降。

三种情况对应的决策边界如\cref{fig:underfit_overfit_plot}所示。欠拟合时，决策边界过于简单，无法有效区分两类样本；拟合良好时，决策边界能够较好地刻画两类样本之间的真实分界，同时忽略少量噪声；过拟合时，决策边界则变得高度复杂，为了正确分类训练集中的每一个样本，甚至会随着噪声点发生剧烈弯折。

\begin{figure}[H]
    \boxedimage[width=0.95\textwidth]{figures/part_01/p1_ch04/underfit_overfit_plot.pdf}
    \caption{欠拟合、拟合良好、过拟合的决策边界。散点为两类训练样本，黑线为模型学到的决策边界。}
    \label{fig:underfit_overfit_plot}
\end{figure}

从训练过程来看，三种情况对应的损失曲线如\cref{fig:underfit_overfit_loss}所示。欠拟合时，训练损失和测试损失都较高；拟合良好时，两者都较低且差距较小；而在过拟合情况下，训练损失仍在持续下降，测试损失却在下降一段时间后开始回升，二者之间的差距不断扩大，这种差距称为\textbf{泛化间隙（Generalization Gap）}，也是模型泛化能力下降的一个重要表现。

\begin{figure}[H]
    \boxedimage[width=0.5\textwidth]{figures/part_01/p1_ch04/underfit_overfit_loss.pdf}
    \caption{欠拟合、拟合良好、过拟合的损失曲线。横轴为模型容量，纵轴为损失；红虚线与蓝实线分别表示训练损失与泛化损失。}
    \label{fig:underfit_overfit_loss}
\end{figure}

对于欠拟合，可以通过增加训练轮数、增大模型容量或者引入更有效的特征等方式提升模型的拟合能力。相比之下，过拟合在深度学习中更加常见：现代神经网络通常具有较大的参数量和很强的拟合能力，模型很容易进一步记忆训练数据中的细节与噪声。因此，我们将重点讨论如何缓解过拟合。

\subsection{L1与L2正则化}

当模型过度拟合训练集，甚至连噪声和偶然规律都试图记住时，往往会变得过于复杂，并产生一些过大的权重。例如，某些权重变得非常大，而另一些权重则接近于$0$。模型对少数特征过度依赖，输入稍有变化，预测结果就可能发生较大的波动。

那么，如何限制这种过度拟合？既然过拟合可能表现为部分权重过大，一个直接的思路就是：在原有损失函数后面增加一个\textbf{惩罚项（Penalty Term）}，专门限制权重的规模。这样，当模型试图通过增大权重来拟合训练集中的细节时，惩罚项也会随之增大，迫使模型在“拟合训练数据”和“保持模型简单”之间寻找平衡。于是，优化目标变为：
$$
\mathcal{L}_{\text{new}}=\mathcal{L}_{\text{old}}+\lambda\cdot\Omega(w)
$$
其中，$\Omega(w)$用于衡量权重的整体大小，$\lambda$用于控制正则化的强弱。那么，$\Omega(w)$应该怎么设计？最直接的方式有两种：第一种是对所有权重取绝对值求和，这就是\textbf{L1正则化}\footnote{L1正则化由Robert Tibshirani于1996年提出，论文为\href{https://doi.org/10.1111/j.2517-6161.1996.tb02080.x}{\textit{Regression Shrinkage and Selection Via the Lasso}}。}：
$$
\mathcal{L}_{\text{new}}=\mathcal{L}_{\text{old}}+\lambda\sum_i|w_i|
$$
第二种是对所有权重取平方求和，这就是\textbf{L2正则化}\footnote{L2正则化由Arthur Hoerl和Robert Kennard于1970年提出，论文为\href{https://doi.org/10.1080/00401706.1970.10488635}{\textit{Ridge Regression: Biased Estimation for Nonorthogonal Problems}}。}：
$$
\mathcal{L}_{\text{new}}=\mathcal{L}_{\text{old}}+\lambda\sum_iw_i^2
$$

虽然两种方法的形式非常相似，但它们对权重的约束方式并不相同。L2正则化对大权重的惩罚增长得更快：权重从$1$增大到$10$时，平方惩罚会从$1$增加到$100$，因此它会更强地抑制极端大的权重，使权重整体保持在较小的范围内。L1正则化则对权重施加近似恒定的惩罚，更像是一股持续将权重拉向$0$的力。

那为什么L2通常只是让权重变得越来越小，而L1能够让一些权重恰好变成$0$？从参数更新的角度来看，L2正则化产生的惩罚梯度为$2\lambda w$，它与当前权重的大小成正比。权重越接近$0$，这股拉回$0$的力量就越弱，因此权重通常会逐渐接近$0$，但不会轻易精确变为$0$。L1则不同，在$w\neq0$时，其惩罚梯度的大小始终为$\lambda$。因此，当某个权重已经很小，而原始损失提供的梯度又不足以抵消这股持续的拉力时，该权重就可能被直接推到$0$，从而形成稀疏的参数。

\begin{figure}[H]
    \boxedimage[width=0.8\textwidth]{figures/part_01/p1_ch04/l1_l2.pdf}
    \caption{L1与L2正则化的几何示意。在权重平面$(w_1, w_2)$上，虚线椭圆为原始损失的等高线，实线为正则化约束区域（左为L1的菱形、右为L2的圆形）；黑点为二者接触处，即正则化后的最优解，左图落在坐标轴上、右图落在坐标轴外。}
    \label{fig:l1_l2}
\end{figure}

这种差异从几何角度看会更加直观。如\cref{fig:l1_l2}所示，加入正则化后，最优解可以看作原始损失函数等高线与正则化约束区域的接触位置。L1正则化对应的菱形具有明显的尖角，而这些尖角恰好落在坐标轴上，因此最优解更容易落在坐标轴上，使某些权重变为$0$；L2正则化对应的圆形边界则没有尖角，最优解通常位于坐标轴之外，因此权重会整体变小，但一般不会直接变为$0$。

而“恰好为$0$”与“只是很小”的区别，正是L1能够进行\textbf{特征选择}的原因之一。在简单的线性模型中，每个权重都对应一个输入特征。当某个权重变为$0$时，对应特征对模型输出的直接贡献也随之消失，因此L1正则化在抑制过拟合的同时，还能够得到更加稀疏的模型。不过，正则化强度$\lambda$也不能设置得过大，否则权重会被过度压缩，模型的表达能力下降，最终反而导致欠拟合。由此可见，L1和L2都通过限制权重规模来抑制过拟合，但作用方式并不相同：L1倾向于产生稀疏权重，可以进行特征选择；L2则倾向于让权重整体变小，使模型更加平滑。

\begin{purpleBox}{权重衰减（Weight Decay）}
在标准SGD下，L2正则化与权重衰减是等价的。加入L2正则化后，损失函数对权重的梯度变为：
$$
\frac{\partial\mathcal{L}_{\text{new}}}{\partial w} = \frac{\partial\mathcal{L}_{\text{old}}}{\partial w} + 2\lambda w
$$
将其代入SGD的参数更新公式，可以得到：
$$
w\leftarrow(1-2\eta\lambda)w-\eta\frac{\partial\mathcal{L}_{\text{old}}}{\partial w}.
$$
可以看到，每次更新时，权重都会先乘上一个小于$1$的系数，再进行正常的梯度更新，因此会在训练过程中持续向$0$衰减，这就是“权重衰减”名称的由来。

不过，需要注意的是，L2正则化与权重衰减的等价关系只在标准SGD下成立。当使用Adam等自适应优化器时，二者不再等价。因此，在现代深度学习实践中，通常将L2正则化与权重衰减看作两个不同的概念。
\end{purpleBox}

\subsection{Dropout}

L1和L2正则化主要通过约束权重的大小来抑制过拟合，但深层网络还可能存在另一种更隐蔽的问题：不同神经元之间形成过强的依赖，多个神经元长期配合完成某一项任务，这种现象称为\textbf{协同适应（Co-adaptation）}。一旦输入发生变化，这种高度依赖的组合就可能失去作用，从而影响模型的泛化能力。

那么，能不能在训练过程中随机关闭一部分神经元，让网络无法长期依赖固定的神经元组合？这就是\textbf{Dropout}的核心思想\footnote{Dropout的思想最早由Geoffrey Hinton等人于2012年在预印本\href{https://arxiv.org/abs/1207.0580}{\textit{Improving neural networks by preventing co-adaptation of feature detectors}}中提出，随后同一团队以Nitish Srivastava为第一作者将其整理发表于\textit{Journal of Machine Learning Research}，标题为\href{https://jmlr.org/papers/v15/srivastava14a.html}{\textit{Dropout: A Simple Way to Prevent Neural Networks from Overfitting}}。}：在每次训练的前向传播中，随机选取一部分神经元，将它们的输出置为$0$。这些神经元只在当前迭代中被暂时关闭，不参与前向计算，也不会产生反向梯度，但在下一次迭代中仍有可能重新参与训练。通过不断改变被关闭的神经元，网络需要在不同的神经元组合下完成学习，从而减少对特定神经元或路径的依赖。

\begin{figure}[H]
    \boxedimage[width=0.95\textwidth]{figures/part_01/p1_ch04/dropout.pdf}
    \caption{Dropout示意图。左图为完整网络，红色高亮路径$x_2 \to h_2^{(1)} \to h_1^{(2)} \to y_1$表示模型可能过度依赖的特定路径；右图为一次Dropout采样得到的子网络：带红叉的神经元在本次迭代中被随机关闭，剩余神经元重新连接以完成前向计算。}
    \label{fig:dropout}
\end{figure}

如\cref{fig:dropout}所示，如果不进行任何干预，网络可能逐渐形成对某些特定神经元或路径的依赖。例如，模型可能主要依靠$x_2\to h_2^{(1)}\to h_2^{(2)}\to y_1$这条路径完成预测，而其他神经元则没有充分学习。Dropout会在训练过程中随机切断部分路径，使这种固定的依赖关系难以长期维持。每次迭代被关闭的神经元都可能不同，剩余神经元必须共同完成当前任务，从而促使网络学习更加分散、更加稳健的特征表示。从另一个角度来看，每次随机关闭一部分神经元，相当于训练了一个不同的子网络；而测试时使用完整网络，可以看作综合了大量不同子网络的学习结果。因此，Dropout具有一定的集成学习效果，能够有效降低模型对训练数据的过度依赖，提高泛化能力。

\begin{purpleBox}{反向Dropout（Inverted Dropout）}
在实际实现中，训练时随机关闭一部分神经元，会使下一层接收到的信号整体变弱。为了保持训练阶段和推理阶段的信号尺度一致，通常会采用\textbf{反向Dropout（Inverted Dropout）}：训练时，对保留下来的神经元输出除以保留概率$p$。

例如，当保留概率$p=0.5$时，每次训练只有一半神经元被保留。假设某个神经元的原始输出为$3$，保留后将其缩放为
$$
\frac{3}{0.5}=6
$$
这样，虽然当前只有一半的神经元参与计算，但通过缩放可以使输出的期望值与不使用Dropout时保持一致，从而避免信号整体变弱。而测试时不再随机关闭神经元，也无需额外缩放，直接使用完整网络进行预测。
\end{purpleBox}

\subsection{Early Stopping}

无论是L1/L2正则化还是Dropout，都是从模型目标函数或网络结构入手来抑制过拟合。那么，有没有一种方法不改变模型本身，只回答一个问题：什么时候应该停止训练？我们可以观察\cref{fig:underfit_overfit_loss}中训练集损失和验证集损失的变化。训练初期，模型不断学习数据中的规律，两条曲线通常会同时下降。但训练进行到一定阶段后，二者可能开始出现分歧：训练损失仍在下降，而验证集损失却不再下降，甚至开始回升。这通常意味着模型开始逐渐拟合训练集中的噪声，过拟合正在发生。

既然如此，一个直接的做法就是在训练过程中持续监控验证集损失。一旦验证集损失不再改善，就停止训练，并保存验证集表现最好的那一组模型参数。这就是\textbf{早停机制（Early Stopping）}。那么，为什么提前停止训练也能起到正则化的作用？训练开始时，模型通常从较小的随机参数出发；随着训练不断进行，模型的拟合能力逐渐增强，也越来越容易记忆训练数据中的细节。早停相当于在模型开始过度拟合之前及时停止训练，从而限制模型继续学习训练集中的噪声和偶然规律。

在实际训练中，为了避免验证集损失的偶然波动导致模型过早停止，通常会设置一个\textbf{耐心值（Patience）}。只有当验证集损失连续若干个Epoch没有改善时，才真正停止训练。例如，将Patience设置为10，就表示连续10个Epoch都没有取得更好的验证结果后才停止。因此，早停不需要修改模型结构，只需监控验证集表现，就能在合适的时机停止训练，是深度学习中一种简单而常用的正则化方法。

\begin{quote}
    到这里，我们已经掌握了神经网络的基本结构、反向传播、优化算法以及训练与泛化等基础知识。接下来，我们来看看神经网络是如何处理图像的。
\end{quote}